% =====================================================================================================
% inner_speech_section.tex — Inner speech, mind-wandering and active inference
%
% DRAFT (2026-11-18). A proposal section: it describes an architecture that is NOT yet implemented,
% grounded in the existing language_agent (robocomp/active_inference). Include from main.tex with
%     % =====================================================================================================
% inner_speech_section.tex — Inner speech, mind-wandering and active inference
%
% DRAFT (2026-11-18). A proposal section: it describes an architecture that is NOT yet implemented,
% grounded in the existing language_agent (robocomp/active_inference). Include from main.tex with
%     % =====================================================================================================
% inner_speech_section.tex — Inner speech, mind-wandering and active inference
%
% DRAFT (2026-11-18). A proposal section: it describes an architecture that is NOT yet implemented,
% grounded in the existing language_agent (robocomp/active_inference). Include from main.tex with
%     % =====================================================================================================
% inner_speech_section.tex — Inner speech, mind-wandering and active inference
%
% DRAFT (2026-11-18). A proposal section: it describes an architecture that is NOT yet implemented,
% grounded in the existing language_agent (robocomp/active_inference). Include from main.tex with
%     \input{inner_speech_section}
% and add inner_speech.bib as a second resource:  \addbibresource{inner_speech.bib}
% All keys are new (none collides with the keys main.tex already cites).
% Packages used: amsmath, amssymb, enumitem (already loaded by main.tex) and array, which main.tex does not
% load yet: add \usepackage{array} to its preamble.
% =====================================================================================================

\section{Inner Speech as Epistemic Action}
\label{sec:inner_speech}

The robot described so far uses language in one direction only: a person speaks, and a language agent turns
the request into commands for the controller, grounded in a textual digest of the belief graph. This section
proposes to turn that capability inward. We argue that a language model embedded in an Active Inference
architecture can play the role that inner speech plays in human cognition: a serial, compositional medium in
which the agent poses questions to itself, generates hypotheses, rehearses plans and replays its past, and that
this role follows naturally from the epistemic component of expected free energy. The proposal is stated
together with the constraints that keep such a voice from becoming a source of confabulated belief, and with the
experiments that would falsify it.

\subsection{Motivation: language for thought before language for talk}
\label{sec:inner_speech:motivation}

Two traditions disagree on the origin of inner speech but agree on its function. On the internalist view,
human language evolved first as an internal system for structuring thought, and its use for communication is a
later externalisation \cite{berwick2016why}. On the developmental view, inner speech is social speech that the
child gradually internalises, abbreviates and turns into a tool for self-regulation \cite{vygotsky1986thought}.
Contemporary reviews treat inner speech as functionally involved in planning, problem solving, self-monitoring and
the maintenance of goals in working memory \cite{aldersonday2015inner}, and language more generally as a
cognitive technology that lets a learner manipulate discrete, recombinable tokens rather than raw percepts
\cite{clark1998magic}.

Large language models give a striking instance of the developmental route. They are trained exclusively on
communicative text, yet producing intermediate reasoning in that same medium before answering --- chain-of-thought
--- measurably improves their performance \cite{wei2022chain}. A system built for talk becomes better at thinking
by talking to itself. Our language agent sits at the same point of its own development: it was built to converse
with a person, and the question we ask is what it contributes once its speech is addressed to the robot itself.
Section~\ref{sec:inner_speech:why_lm} returns to what this origin implies for the internalist view.

\subsection{Inner questions as epistemic actions}
\label{sec:inner_speech:questions}

In Active Inference, policies are scored by expected free energy, whose epistemic term rewards the expected
reduction of uncertainty about hidden states \cite{friston2017process, parr2022active}. A question is the
linguistic form of exactly this quantity: it names a hidden variable whose value is uncertain and whose resolution
is expected to matter. We therefore model an inner question as a persistent object
\begin{equation}
    \mathcal{Q} = \big\langle\, e,\; \mathcal{H},\; q(h),\; \mathcal{P},\; b,\; s \,\big\rangle ,
    \label{eq:inner_question}
\end{equation}
where $e$ is the triggering event (for example an obstacle no concept agent explains), $\mathcal{H}$ a finite set
of named hypotheses (``a chair'', ``a box'', ``unknown''), $q(h)$ the current posterior over them, $\mathcal{P}$
the set of pending predictions, $b$ a budget in time and motion, and $s \in \{\text{open},\text{resolved},
\text{abandoned}\}$ its status. Questions form a stack in the sense of a \emph{question under discussion}
\cite{roberts2012information}: ``is it a chair?'' is a sub-question whose answer partially resolves the parent
``what is this?''. The question, not the language model's context window, carries the state of the reasoning from
one step to the next.

\paragraph{One step of the loop.}
Given $\mathcal{Q}$, the language model proposes a test: an epistemic action $a$ (turn towards the anomaly,
approach it, look with the front camera) together with the outcome each hypothesis predicts. The value of the test
is its expected information gain minus its cost,
\begin{equation}
    V(a) \;=\; \underbrace{H\!\left[q(h)\right] \;-\; \mathbb{E}_{p(o \mid a)}\!\left[ H\!\left[q(h \mid o, a)\right]
    \right]}_{I(h;\,o \mid a)} \;-\; c(a),
    \label{eq:inner_value}
\end{equation}
which is the negative of the epistemic part of expected free energy with an explicit action cost. The test is
executed only if $V(a) > 0$ and the action passes the same deterministic gate that already validates every
command the language agent issues. When the observation $o$ arrives, the posterior is updated by the concept
agents' own likelihoods,
\begin{equation}
    q(h \mid o, a) \;\propto\; p(o \mid h, a)\, q(h),
    \label{eq:inner_update}
\end{equation}
and the loop continues, spawns a sub-question, or terminates when
\begin{equation}
    H\!\left[q(h)\right] < \eta \quad\lor\quad \max_a V(a) \le 0 \quad\lor\quad b \text{ exhausted}.
    \label{eq:inner_stop}
\end{equation}
The stopping rule is what separates reasoning from rumination: a loop with no expected information to gain closes,
possibly with the answer ``unknown'', and its trigger is habituated for a refractory period.

\paragraph{The grounding constraint.}
The central design rule is that \emph{the language model proposes and perception disposes}. Text produced by the
model may introduce hypotheses, predictions and tests; it may never move a posterior. Only the likelihood
$p(o \mid h, a)$ in \eqref{eq:inner_update}, computed by the perceptual agents from sensor data, does. This is the
belief-side counterpart of the rule the language agent already enforces for actions, where every proposed command
is validated in C++ against the belief graph before it reaches the controller. Without it, a self-addressed
language loop is a closed system of text conditioned on text and converges on its own confabulations.

\paragraph{The proposer as a calibrated prior.}
The hypotheses the model names, and the probabilities it attaches to them, encode world knowledge the robot has
not acquired by itself (``chairs come with tables''). We treat this as a prior whose weight is earned. Let
$\pi_{\text{LM}}(h \mid x)$ be the proposer's distribution in context $x$ and $n_x(h)$ the robot's own counts of
resolved questions in that context. The prior entering \eqref{eq:inner_update} is
\begin{equation}
    q_0(h \mid x) \;\propto\; \pi_{\text{LM}}(h \mid x)^{\beta_x}\, \big(\alpha + n_x(h)\big),
    \label{eq:inner_prior}
\end{equation}
a tempered product of the proposer and a Dirichlet over experience, where $\beta_x \in [0,1]$ is set from the
proposer's measured calibration in that context (for instance from its Brier score on past questions). Early in
the robot's life the proposer dominates; as counts accumulate, the robot's own experience replaces borrowed
knowledge. The same conjugate form already governs the language agent's decision to ask a person or to assume
that they moved an object, where a Beta posterior over the person's past answers is compared against the costs of
asking and of being wrong.

\subsection{Attention, automaticity and the wandering mind}
\label{sec:inner_speech:modes}

A human driver on an easy road steers without attending to it; the mind is free to plan, imagine or recall, and is
recaptured the instant something unexpected happens \cite{smallwood2015science, christoff2016mind}. In Active
Inference this is a statement about precision: while prediction errors are small and the habitual policy is
reliable, attention --- the precision assigned to sensory prediction errors --- can be withdrawn from the task; a
surprise raises it again \cite{feldman2010attention}. The arbitration between habitual and deliberative control by
their relative uncertainty \cite{daw2005uncertainty} has the same structure.

The architecture already separates the two systems. The controller and the concept agents form an automatic layer
that drives, avoids obstacles and executes missions without the language model; the language agent is the
deliberative voice. We propose to make the allocation of the voice explicit through a scalar \emph{demand} signal,
computed on every digest cycle without the language model,
\begin{equation}
    D_t \;=\; \sum_{k} \pi_k\, \varepsilon_{k,t}^{2} \;+\; \lambda\, \mathbb{1}\!\left[\text{person active}_t\right],
    \label{eq:inner_demand}
\end{equation}
where $\varepsilon_{k,t}$ are normalised prediction errors reported by the agents (unexplained residual occupancy,
flips of existence or door beliefs, localisation uncertainty, stalled or failed mission steps, controller refusals)
and $\pi_k$ their precisions. Three modes follow from two thresholds with asymmetric hysteresis:
\begin{equation}
    m_t =
    \begin{cases}
        \textsc{focused}    & D_t > \theta_{\text{hi}},\\[2pt]
        \textsc{wandering}  & D_\tau < \theta_{\text{lo}} \;\; \forall\, \tau \in [t - T, t],\\[2pt]
        \textsc{monitoring} & \text{otherwise.}
    \end{cases}
    \label{eq:inner_modes}
\end{equation}
Entering \textsc{wandering} requires demand to stay low for a dwell time $T$; leaving it requires a single sample
above $\theta_{\text{hi}}$. The robot is slow to drift off and instant to snap back. Safety never depends on this
arbitration: the controller's reflexes act on their own, and the deliberative voice arrives after the reflex, as
awareness follows braking in the driver.

\paragraph{Focused mode} runs the question loop of Section~\ref{sec:inner_speech:questions} on the triggering
event. \paragraph{Monitoring mode} makes no language-model calls; the voice is primed but silent.

\paragraph{Wandering mode.}
Freed from the task, the voice is seeded with one recent experience chosen by salience: an episode that ended
unresolved, a failure, an object that changed, or an upcoming goal. Each step's output, together with an
associative retrieval from episodic memory, seeds the next, so that the chain drifts as spontaneous thought does
\cite{christoff2016mind}. The voice has access only to memory queries and to the architecture's imagination tools,
which simulate observations and outcomes on a copy of the belief graph; it can neither act nor write beliefs. Its
products are of three kinds: counterfactuals over past episodes (``had the door been open, the crossing would have
taken 40\,s less''), hypotheses to be tested later, and prospective plans (``if the door is shut on the way back,
the route through the next room is shorter''). Functionally this is the default network's role in simulating
past and future \cite{buckner2008brain}, and the offline replay and model reduction that Active Inference assigns
to rest and sleep \cite{friston2017curiosity, hobson2012waking}. Replay is prioritised by the expected gain of the
memory it revisits \cite{mattar2018prioritized}.

\paragraph{Source monitoring.}
Humans occasionally misremember imagined events as perceived ones \cite{johnson1993source}. The robot must not.
Everything produced while wandering is stored with an \emph{imagined} provenance, separate from the factual
episodic record, and enters the learning loop only as a hypothesis to be tested, never as evidence.

\subsection{What the return to reality gives the automatic layer}
\label{sec:inner_speech:return}

The transition from \textsc{wandering} to \textsc{focused} cancels the chain in flight (its state is kept, so that
it may resume if demand subsides) and gives the full voice to the surprise. It serves the automatic layer at three
timescales.
\begin{enumerate}
    \item \textbf{Decision.} A surprise means the automatic layer has left its domain of competence, where its only
    responses are to retry, stop or fail. The voice adds re-interpretation of the scene, arbitration between goals,
    re-planning through remembered topology, and recourse to a person.
    \item \textbf{Attention.} The voice can raise precision where it matters: orienting sensors, increasing a
    detector's rate in a region, or passing a hypothesis to the concept agents as a spatial prior. Such a prior has
    bounded weight; it may accelerate the recognition of an object that is present but may not create one that is
    not, preserving the grounding constraint.
    \item \textbf{Teaching.} Each resolved surprise is a labelled example the automatic layer could not have
    produced alone: a residual later identified as a chair is a false negative of the chair detector; a door that
    is repeatedly shut in the evening becomes a prior; a recovery that repeatedly succeeds becomes a branch of the
    controller's behaviour tree. The domain of competence grows, and the situations that require the voice become
    rarer, which is the classical progression from effortful to autonomous skill \cite{fitts1967human}.
\end{enumerate}
Plans prepared while wandering act as a cache: when the anticipated situation occurs, the focused voice verifies a
plan instead of composing one. The return ends with an explicit release --- spatial priors cleared, safety margins
restored, the automatic layer confirmed back within its competence --- without which every surprise would leave
the robot permanently vigilant.

\subsection{Relation to chain-of-thought reasoning}
\label{sec:inner_speech:cot}

Chain-of-thought is inner speech in the functional sense, but of a specific kind, and the comparison sharpens the
design (Table~\ref{tab:inner_cot}). Plain chain-of-thought is convergent and ungrounded: tokens condition on
tokens until an answer is produced, and the trace is discarded. Interleaving reasoning with action and observation
grounds it \cite{yao2023react}; our focused loop is of this kind, with the additional property that the
observation, through \eqref{eq:inner_update}, rather than the next token, determines where the reasoning goes.
The wandering voice is closer to plain chain-of-thought in its isolation from perception, but divergent rather
than convergent, interruptible, and persistent in its products.

\begin{table}[ht]
    \centering
    \small
    \begin{tabular}{|>{\raggedright\arraybackslash}p{0.19\linewidth}|>{\raggedright\arraybackslash}p{0.22\linewidth}|>{\raggedright\arraybackslash}p{0.24\linewidth}|>{\raggedright\arraybackslash}p{0.24\linewidth}|}
        \hline
                                & \textbf{Chain-of-thought} & \textbf{Focused voice} & \textbf{Wandering voice} \\ \hline
        Goal                    & answer a given question   & resolve a given surprise & none given; seed by salience \\ \hline
        Dynamics                & convergent                & convergent             & divergent, associative \\ \hline
        Contact with the world  & none                      & thought--action--observation & none (imagination only) \\ \hline
        Persistence             & discarded                 & question objects, beliefs, lessons & imagined store, hypotheses \\ \hline
        Interruptible           & no                        & yes                    & yes, resumable \\ \hline
        Latency budget          & n/a                       & $\sim$1\,s, no extended thinking & unbounded, extended thinking \\ \hline
    \end{tabular}
    \caption{Chain-of-thought compared with the two modes of the proposed inner voice.}
    \label{tab:inner_cot}
\end{table}

Two practical consequences follow. First, \emph{think when there is time, not when there is pressure}. In the
deployed language agent extended thinking is disabled, because a thinking-mode model spent 3\,000--6\,000 tokens
(40--80\,s) before each tool call; that is incompatible with the focused mode but costs nothing in the wandering
mode, where it can be enabled with a large budget and pre-empted at any moment. Slow reasoning is done in advance
and fast reasoning uses its results. Second, chain-of-thought traces are known not always to reflect the
computation that produced the answer \cite{turpin2023language}, as human verbal reports do not always reflect the
causes of behaviour \cite{nisbett1977telling}. The faithful record of the robot's reasoning is therefore the
structured one --- hypotheses, predictions, observations and posterior updates --- not the text. Finally, wandering
chains whose hypotheses were later confirmed by perception form a corpus of reasoning verified against the world,
which could be used to fine-tune the local model in the manner of self-taught reasoning \cite{zelikman2022star},
with perception, rather than a known answer, as the filter.

\subsection{Consequences for belief and learning}
\label{sec:inner_speech:learning}

The loop changes how beliefs evolve in two ways. Within an episode, uncertainty is reduced where surprise
occurred rather than where the sensors happened to point, which is efficient but leaves the rest of the map less
examined and makes the robot's own data a biased sample; a decaying novelty term \cite{schwartenbeck2019curiosity}
and importance weighting of learned statistics are required to correct this. Across episodes, every closed
question is a labelled example, and learning proceeds at four timescales (Table~\ref{tab:inner_learning}).

\begin{table}[ht]
    \centering
    \small
    \begin{tabular}{|>{\raggedright\arraybackslash}p{0.14\linewidth}|>{\raggedright\arraybackslash}p{0.37\linewidth}|>{\raggedright\arraybackslash}p{0.37\linewidth}|}
        \hline
        \textbf{Timescale} & \textbf{What changes} & \textbf{Where, in this architecture} \\ \hline
        Seconds   & posterior over the hypotheses of one question         & the question object \eqref{eq:inner_question} \\ \hline
        Days      & priors: where objects are, which doors are shut, what moves & long-term object and topological memory, as Dirichlet counts \\ \hline
        Weeks     & likelihoods: detector and concept models              & false negatives and positives labelled by resolved questions \\ \hline
        Slow      & structure: new categories                             & recurrent, mutually similar ``unknown'' resolutions named by the voice and proposed as a new concept \\ \hline
    \end{tabular}
    \caption{Learning driven by the inner-question loop.}
    \label{tab:inner_learning}
\end{table}

The last row is where the hypothesis of a language for thought is most directly engaged: a cluster of anomalies
that resists every existing concept becomes learnable once it is named. The loop also learns about itself: the
salience weights $\pi_k$ in \eqref{eq:inner_demand} are adjusted by the information each trigger actually yielded,
asking a person is weighed against looking, and loops that repeatedly resolve the same way are compiled into
reflexive rules, so that the number of language-model calls per hour should \emph{decrease} as competence grows.

\subsection{Why a language model: generation versus evaluation}
\label{sec:inner_speech:why_lm}

None of the above requires a person, so the question arises of what a language model contributes to an
architecture that already performs inference and action selection. Our answer is a division of labour. Expected
free energy evaluates policies and Bayesian updating weighs hypotheses, but both operate over a space that has
already been enumerated: each concept agent carries a fixed generative model, policies must be listed before they
can be scored, and Bayesian model reduction can prune structure but not invent it \cite{friston2017curiosity}.
Whatever falls outside the given models remains an unexplained residual. A language model is weak where this
machinery is strong --- its confidences are not calibrated posteriors and it has no access to the sensors --- and
strong where it is weak: it is a cheap, broad generator of hypotheses, explanations and plans. We therefore use it
only as a \emph{proposal distribution}, in a generate-and-test loop whose test is active inference:
\begin{itemize}
    \item \emph{Hypotheses outside the built-in models.} A residual blob can be proposed to be a bag, a toppled
    chair or a robot vacuum, none of which any concept agent represents; the question loop then tests them.
    \item \emph{Amortised policy search.} Instead of enumerating multi-step policies, the model writes a few
    candidate missions --- small programs in the mission language --- which the gate validates and expected free
    energy scores.
    \item \emph{Relational binding across agents.} The agents' representations (signed distance fields, occupancy
    grids, existence log-odds, episodic records) are mutually incompatible; a sentence such as ``the chair that was
    not next to the table in the room behind door~3 yesterday'' binds them in one query no single agent could pose.
    \item \emph{Explanation from few examples.} After a single surprise the model can propose which variable
    mattered (``the door is shut because it is evening''), turning one episode into a testable generalisation.
    \item \emph{Abstraction and naming.} Episodes are compressed into reusable lessons, and recurrent unknowns are
    named, which is what makes them candidates for a new concept (Table~\ref{tab:inner_learning}).
\end{itemize}
Language, rather than an arbitrary learned embedding, matters for three reasons: it is discrete and compositional,
so new hypotheses and plans are assembled from known parts; it is inspectable, so that deterministic gates can parse
and check it and people can audit it; and, once the robot thinks in it, communication with a person comes at no
further cost.

\paragraph{Where the knowledge comes from.}
This raises a difficulty for the motivating thesis. If the inner voice precedes the outer one, the common-sense
knowledge on which the proposer's value rests cannot have arrived through communication --- yet a language model
acquires all of it from human text, that is, from other people's outer voices. The difficulty dissolves once the
\emph{faculty} is distinguished from its \emph{content}. The internalist claim concerns the former: a
computational system for building unbounded hierarchical structures, available for thought before it is used for
talk \cite{hauser2002faculty, berwick2016why}. It makes no claim that what a person knows was acquired without
language. Human knowledge has three sources: core systems for objects, agents, number and space, present before
learning \cite{spelke2007core}; individual experience; and, once outer language exists, cumulative cultural
transmission, which supplies most of what an adult knows \cite{tomasello1999cultural}. Even on the internalist
view, the inner voice of a modern human is filled largely by the outer voices of others.

The architecture has the same three sources, in a different arrangement. The concept agents' hand-built generative
models play the role of \emph{core knowledge}: given before experience and limited to a few domains. The robot's
own counts $n_x(h)$ in \eqref{eq:inner_prior} are \emph{individual experience}. The language model is
\emph{culture}: knowledge distilled from a community's communication and inherited, for this individual, at birth.
The robot thus resembles a child born into a culture more than an early hominin, and its inner voice is Vygotskian
in origin: its format and its content were both learned from outer speech. Two consequences follow. First, the
proposal does not demonstrate the internalist thesis and should not be read as doing so; its claim is functional
--- an inner use of language has epistemic value independently of communication --- and holds whatever the origin of
the language system. Second, the format and the content of the voice can in principle be separated, which turns the
question into an experiment (P7 in Section~\ref{sec:inner_speech:predictions}): a compositional generator over the
robot's own concepts and mission language, trained only on the robot's experience, in the manner of library
learning in program induction \cite{ellis2021dreamcoder}, would supply the format without the culture.

\subsection{Mapping onto the existing architecture}
\label{sec:inner_speech:mapping}

Several components already exist in the language agent. Robot events (mission reports, failures, completions)
already wake the model through a queue, in turns restricted to perceiving, querying memory, stopping and
proposing, never acting. Every action already passes a deterministic gate. The decision to ask or to assume is
already decision-theoretic. Unexplained occupancy is already isolated by a residual agent that acts as the null
hypothesis of the scene, and episodic and long-term object memories are already joined read-only. Missing are the
persistent question objects \eqref{eq:inner_question}, a prediction evaluator implementing
\eqref{eq:inner_update} outside the language model, the demand signal and mode arbitration
\eqref{eq:inner_demand}--\eqref{eq:inner_modes}, the imagined store with its provenance, and the channel through
which bounded spatial priors reach the concept agents.

\subsection{Testable predictions}
\label{sec:inner_speech:predictions}

Consistent with the evaluation principle adopted throughout --- policies compared against each other inside one
closed loop, with body, sensors and environment fixed --- the proposal makes seven falsifiable predictions, each
tested in simulation by moving known objects and introducing unfamiliar ones:
\begin{enumerate}[label=\textbf{P\arabic*}]
    \item The unexplained residual mass per room decreases faster with the focused loop than without it.
    \item The proposer's calibration on resolved questions improves, and $\beta_x$ in \eqref{eq:inner_prior}
    decreases, as the robot's own counts accumulate.
    \item Detector false-negative rates fall after retraining on question-labelled examples.
    \item Inner questions per hour decrease over repeated exposure to the same environment while answer accuracy
    is maintained.
    \item Plans prepared while wandering reduce the latency from surprise to a validated response, and the time to
    notice a deliberately moved object does not grow despite habituation.
    \item With the language-model proposer replaced by a fixed hypothesis set (the existing concept classes and
    ``unknown'') and template-based missions, the loop resolves fewer questions about objects and failures outside
    the built-in models, and acquires new categories more slowly.
    \item A compositional generator trained only on the robot's own experience (format without culture) recovers
    part of the advantage in P6 within a familiar environment, and less of it the further a situation departs from
    what the robot has seen.
\end{enumerate}
A failure of P1 would indicate that the epistemic value of language-mediated questions does not exceed that of
the existing exploration; a failure of P4 would indicate that the loop does not transfer competence to the
automatic layer, which is its principal claim; and a failure of P6 would indicate that the language model adds
nothing a fixed hypothesis set does not, which would remove the case for including it. P7 separates the
contribution of language as a format for thought from that of the knowledge it carries.

\subsection{Limitations}
\label{sec:inner_speech:limitations}

The demand signal \eqref{eq:inner_demand} sums heterogeneous prediction errors whose precisions must themselves be
calibrated, and a poorly scaled term will either suppress the voice or keep it permanently engaged. A single
language model serves the person, the focused loop and the wandering loop; pre-emption keeps the person first but
bounds how much offline reasoning a busy robot can do. The tempered prior \eqref{eq:inner_prior} is a heuristic
combination rather than a derived posterior. The generator of P7 is substantially harder to build than the
ablation of P6, and may have to be approximated by a small model trained on the robot's logs and mission corpus,
whose pre-training would leak some culture back in. Finally, the analogy with human inner speech is functional, not
mechanistic: the proposal is that the same computational role --- serial, compositional self-questioning in the
service of epistemic value --- is useful for a robot, not that the robot's voice resembles the human one in any
further respect.

% and add inner_speech.bib as a second resource:  \addbibresource{inner_speech.bib}
% All keys are new (none collides with the keys main.tex already cites).
% Packages used: amsmath, amssymb, enumitem (already loaded by main.tex) and array, which main.tex does not
% load yet: add \usepackage{array} to its preamble.
% =====================================================================================================

\section{Inner Speech as Epistemic Action}
\label{sec:inner_speech}

The robot described so far uses language in one direction only: a person speaks, and a language agent turns
the request into commands for the controller, grounded in a textual digest of the belief graph. This section
proposes to turn that capability inward. We argue that a language model embedded in an Active Inference
architecture can play the role that inner speech plays in human cognition: a serial, compositional medium in
which the agent poses questions to itself, generates hypotheses, rehearses plans and replays its past, and that
this role follows naturally from the epistemic component of expected free energy. The proposal is stated
together with the constraints that keep such a voice from becoming a source of confabulated belief, and with the
experiments that would falsify it.

\subsection{Motivation: language for thought before language for talk}
\label{sec:inner_speech:motivation}

Two traditions disagree on the origin of inner speech but agree on its function. On the internalist view,
human language evolved first as an internal system for structuring thought, and its use for communication is a
later externalisation \cite{berwick2016why}. On the developmental view, inner speech is social speech that the
child gradually internalises, abbreviates and turns into a tool for self-regulation \cite{vygotsky1986thought}.
Contemporary reviews treat inner speech as functionally involved in planning, problem solving, self-monitoring and
the maintenance of goals in working memory \cite{aldersonday2015inner}, and language more generally as a
cognitive technology that lets a learner manipulate discrete, recombinable tokens rather than raw percepts
\cite{clark1998magic}.

Large language models give a striking instance of the developmental route. They are trained exclusively on
communicative text, yet producing intermediate reasoning in that same medium before answering --- chain-of-thought
--- measurably improves their performance \cite{wei2022chain}. A system built for talk becomes better at thinking
by talking to itself. Our language agent sits at the same point of its own development: it was built to converse
with a person, and the question we ask is what it contributes once its speech is addressed to the robot itself.
Section~\ref{sec:inner_speech:why_lm} returns to what this origin implies for the internalist view.

\subsection{Inner questions as epistemic actions}
\label{sec:inner_speech:questions}

In Active Inference, policies are scored by expected free energy, whose epistemic term rewards the expected
reduction of uncertainty about hidden states \cite{friston2017process, parr2022active}. A question is the
linguistic form of exactly this quantity: it names a hidden variable whose value is uncertain and whose resolution
is expected to matter. We therefore model an inner question as a persistent object
\begin{equation}
    \mathcal{Q} = \big\langle\, e,\; \mathcal{H},\; q(h),\; \mathcal{P},\; b,\; s \,\big\rangle ,
    \label{eq:inner_question}
\end{equation}
where $e$ is the triggering event (for example an obstacle no concept agent explains), $\mathcal{H}$ a finite set
of named hypotheses (``a chair'', ``a box'', ``unknown''), $q(h)$ the current posterior over them, $\mathcal{P}$
the set of pending predictions, $b$ a budget in time and motion, and $s \in \{\text{open},\text{resolved},
\text{abandoned}\}$ its status. Questions form a stack in the sense of a \emph{question under discussion}
\cite{roberts2012information}: ``is it a chair?'' is a sub-question whose answer partially resolves the parent
``what is this?''. The question, not the language model's context window, carries the state of the reasoning from
one step to the next.

\paragraph{One step of the loop.}
Given $\mathcal{Q}$, the language model proposes a test: an epistemic action $a$ (turn towards the anomaly,
approach it, look with the front camera) together with the outcome each hypothesis predicts. The value of the test
is its expected information gain minus its cost,
\begin{equation}
    V(a) \;=\; \underbrace{H\!\left[q(h)\right] \;-\; \mathbb{E}_{p(o \mid a)}\!\left[ H\!\left[q(h \mid o, a)\right]
    \right]}_{I(h;\,o \mid a)} \;-\; c(a),
    \label{eq:inner_value}
\end{equation}
which is the negative of the epistemic part of expected free energy with an explicit action cost. The test is
executed only if $V(a) > 0$ and the action passes the same deterministic gate that already validates every
command the language agent issues. When the observation $o$ arrives, the posterior is updated by the concept
agents' own likelihoods,
\begin{equation}
    q(h \mid o, a) \;\propto\; p(o \mid h, a)\, q(h),
    \label{eq:inner_update}
\end{equation}
and the loop continues, spawns a sub-question, or terminates when
\begin{equation}
    H\!\left[q(h)\right] < \eta \quad\lor\quad \max_a V(a) \le 0 \quad\lor\quad b \text{ exhausted}.
    \label{eq:inner_stop}
\end{equation}
The stopping rule is what separates reasoning from rumination: a loop with no expected information to gain closes,
possibly with the answer ``unknown'', and its trigger is habituated for a refractory period.

\paragraph{The grounding constraint.}
The central design rule is that \emph{the language model proposes and perception disposes}. Text produced by the
model may introduce hypotheses, predictions and tests; it may never move a posterior. Only the likelihood
$p(o \mid h, a)$ in \eqref{eq:inner_update}, computed by the perceptual agents from sensor data, does. This is the
belief-side counterpart of the rule the language agent already enforces for actions, where every proposed command
is validated in C++ against the belief graph before it reaches the controller. Without it, a self-addressed
language loop is a closed system of text conditioned on text and converges on its own confabulations.

\paragraph{The proposer as a calibrated prior.}
The hypotheses the model names, and the probabilities it attaches to them, encode world knowledge the robot has
not acquired by itself (``chairs come with tables''). We treat this as a prior whose weight is earned. Let
$\pi_{\text{LM}}(h \mid x)$ be the proposer's distribution in context $x$ and $n_x(h)$ the robot's own counts of
resolved questions in that context. The prior entering \eqref{eq:inner_update} is
\begin{equation}
    q_0(h \mid x) \;\propto\; \pi_{\text{LM}}(h \mid x)^{\beta_x}\, \big(\alpha + n_x(h)\big),
    \label{eq:inner_prior}
\end{equation}
a tempered product of the proposer and a Dirichlet over experience, where $\beta_x \in [0,1]$ is set from the
proposer's measured calibration in that context (for instance from its Brier score on past questions). Early in
the robot's life the proposer dominates; as counts accumulate, the robot's own experience replaces borrowed
knowledge. The same conjugate form already governs the language agent's decision to ask a person or to assume
that they moved an object, where a Beta posterior over the person's past answers is compared against the costs of
asking and of being wrong.

\subsection{Attention, automaticity and the wandering mind}
\label{sec:inner_speech:modes}

A human driver on an easy road steers without attending to it; the mind is free to plan, imagine or recall, and is
recaptured the instant something unexpected happens \cite{smallwood2015science, christoff2016mind}. In Active
Inference this is a statement about precision: while prediction errors are small and the habitual policy is
reliable, attention --- the precision assigned to sensory prediction errors --- can be withdrawn from the task; a
surprise raises it again \cite{feldman2010attention}. The arbitration between habitual and deliberative control by
their relative uncertainty \cite{daw2005uncertainty} has the same structure.

The architecture already separates the two systems. The controller and the concept agents form an automatic layer
that drives, avoids obstacles and executes missions without the language model; the language agent is the
deliberative voice. We propose to make the allocation of the voice explicit through a scalar \emph{demand} signal,
computed on every digest cycle without the language model,
\begin{equation}
    D_t \;=\; \sum_{k} \pi_k\, \varepsilon_{k,t}^{2} \;+\; \lambda\, \mathbb{1}\!\left[\text{person active}_t\right],
    \label{eq:inner_demand}
\end{equation}
where $\varepsilon_{k,t}$ are normalised prediction errors reported by the agents (unexplained residual occupancy,
flips of existence or door beliefs, localisation uncertainty, stalled or failed mission steps, controller refusals)
and $\pi_k$ their precisions. Three modes follow from two thresholds with asymmetric hysteresis:
\begin{equation}
    m_t =
    \begin{cases}
        \textsc{focused}    & D_t > \theta_{\text{hi}},\\[2pt]
        \textsc{wandering}  & D_\tau < \theta_{\text{lo}} \;\; \forall\, \tau \in [t - T, t],\\[2pt]
        \textsc{monitoring} & \text{otherwise.}
    \end{cases}
    \label{eq:inner_modes}
\end{equation}
Entering \textsc{wandering} requires demand to stay low for a dwell time $T$; leaving it requires a single sample
above $\theta_{\text{hi}}$. The robot is slow to drift off and instant to snap back. Safety never depends on this
arbitration: the controller's reflexes act on their own, and the deliberative voice arrives after the reflex, as
awareness follows braking in the driver.

\paragraph{Focused mode} runs the question loop of Section~\ref{sec:inner_speech:questions} on the triggering
event. \paragraph{Monitoring mode} makes no language-model calls; the voice is primed but silent.

\paragraph{Wandering mode.}
Freed from the task, the voice is seeded with one recent experience chosen by salience: an episode that ended
unresolved, a failure, an object that changed, or an upcoming goal. Each step's output, together with an
associative retrieval from episodic memory, seeds the next, so that the chain drifts as spontaneous thought does
\cite{christoff2016mind}. The voice has access only to memory queries and to the architecture's imagination tools,
which simulate observations and outcomes on a copy of the belief graph; it can neither act nor write beliefs. Its
products are of three kinds: counterfactuals over past episodes (``had the door been open, the crossing would have
taken 40\,s less''), hypotheses to be tested later, and prospective plans (``if the door is shut on the way back,
the route through the next room is shorter''). Functionally this is the default network's role in simulating
past and future \cite{buckner2008brain}, and the offline replay and model reduction that Active Inference assigns
to rest and sleep \cite{friston2017curiosity, hobson2012waking}. Replay is prioritised by the expected gain of the
memory it revisits \cite{mattar2018prioritized}.

\paragraph{Source monitoring.}
Humans occasionally misremember imagined events as perceived ones \cite{johnson1993source}. The robot must not.
Everything produced while wandering is stored with an \emph{imagined} provenance, separate from the factual
episodic record, and enters the learning loop only as a hypothesis to be tested, never as evidence.

\subsection{What the return to reality gives the automatic layer}
\label{sec:inner_speech:return}

The transition from \textsc{wandering} to \textsc{focused} cancels the chain in flight (its state is kept, so that
it may resume if demand subsides) and gives the full voice to the surprise. It serves the automatic layer at three
timescales.
\begin{enumerate}
    \item \textbf{Decision.} A surprise means the automatic layer has left its domain of competence, where its only
    responses are to retry, stop or fail. The voice adds re-interpretation of the scene, arbitration between goals,
    re-planning through remembered topology, and recourse to a person.
    \item \textbf{Attention.} The voice can raise precision where it matters: orienting sensors, increasing a
    detector's rate in a region, or passing a hypothesis to the concept agents as a spatial prior. Such a prior has
    bounded weight; it may accelerate the recognition of an object that is present but may not create one that is
    not, preserving the grounding constraint.
    \item \textbf{Teaching.} Each resolved surprise is a labelled example the automatic layer could not have
    produced alone: a residual later identified as a chair is a false negative of the chair detector; a door that
    is repeatedly shut in the evening becomes a prior; a recovery that repeatedly succeeds becomes a branch of the
    controller's behaviour tree. The domain of competence grows, and the situations that require the voice become
    rarer, which is the classical progression from effortful to autonomous skill \cite{fitts1967human}.
\end{enumerate}
Plans prepared while wandering act as a cache: when the anticipated situation occurs, the focused voice verifies a
plan instead of composing one. The return ends with an explicit release --- spatial priors cleared, safety margins
restored, the automatic layer confirmed back within its competence --- without which every surprise would leave
the robot permanently vigilant.

\subsection{Relation to chain-of-thought reasoning}
\label{sec:inner_speech:cot}

Chain-of-thought is inner speech in the functional sense, but of a specific kind, and the comparison sharpens the
design (Table~\ref{tab:inner_cot}). Plain chain-of-thought is convergent and ungrounded: tokens condition on
tokens until an answer is produced, and the trace is discarded. Interleaving reasoning with action and observation
grounds it \cite{yao2023react}; our focused loop is of this kind, with the additional property that the
observation, through \eqref{eq:inner_update}, rather than the next token, determines where the reasoning goes.
The wandering voice is closer to plain chain-of-thought in its isolation from perception, but divergent rather
than convergent, interruptible, and persistent in its products.

\begin{table}[ht]
    \centering
    \small
    \begin{tabular}{|>{\raggedright\arraybackslash}p{0.19\linewidth}|>{\raggedright\arraybackslash}p{0.22\linewidth}|>{\raggedright\arraybackslash}p{0.24\linewidth}|>{\raggedright\arraybackslash}p{0.24\linewidth}|}
        \hline
                                & \textbf{Chain-of-thought} & \textbf{Focused voice} & \textbf{Wandering voice} \\ \hline
        Goal                    & answer a given question   & resolve a given surprise & none given; seed by salience \\ \hline
        Dynamics                & convergent                & convergent             & divergent, associative \\ \hline
        Contact with the world  & none                      & thought--action--observation & none (imagination only) \\ \hline
        Persistence             & discarded                 & question objects, beliefs, lessons & imagined store, hypotheses \\ \hline
        Interruptible           & no                        & yes                    & yes, resumable \\ \hline
        Latency budget          & n/a                       & $\sim$1\,s, no extended thinking & unbounded, extended thinking \\ \hline
    \end{tabular}
    \caption{Chain-of-thought compared with the two modes of the proposed inner voice.}
    \label{tab:inner_cot}
\end{table}

Two practical consequences follow. First, \emph{think when there is time, not when there is pressure}. In the
deployed language agent extended thinking is disabled, because a thinking-mode model spent 3\,000--6\,000 tokens
(40--80\,s) before each tool call; that is incompatible with the focused mode but costs nothing in the wandering
mode, where it can be enabled with a large budget and pre-empted at any moment. Slow reasoning is done in advance
and fast reasoning uses its results. Second, chain-of-thought traces are known not always to reflect the
computation that produced the answer \cite{turpin2023language}, as human verbal reports do not always reflect the
causes of behaviour \cite{nisbett1977telling}. The faithful record of the robot's reasoning is therefore the
structured one --- hypotheses, predictions, observations and posterior updates --- not the text. Finally, wandering
chains whose hypotheses were later confirmed by perception form a corpus of reasoning verified against the world,
which could be used to fine-tune the local model in the manner of self-taught reasoning \cite{zelikman2022star},
with perception, rather than a known answer, as the filter.

\subsection{Consequences for belief and learning}
\label{sec:inner_speech:learning}

The loop changes how beliefs evolve in two ways. Within an episode, uncertainty is reduced where surprise
occurred rather than where the sensors happened to point, which is efficient but leaves the rest of the map less
examined and makes the robot's own data a biased sample; a decaying novelty term \cite{schwartenbeck2019curiosity}
and importance weighting of learned statistics are required to correct this. Across episodes, every closed
question is a labelled example, and learning proceeds at four timescales (Table~\ref{tab:inner_learning}).

\begin{table}[ht]
    \centering
    \small
    \begin{tabular}{|>{\raggedright\arraybackslash}p{0.14\linewidth}|>{\raggedright\arraybackslash}p{0.37\linewidth}|>{\raggedright\arraybackslash}p{0.37\linewidth}|}
        \hline
        \textbf{Timescale} & \textbf{What changes} & \textbf{Where, in this architecture} \\ \hline
        Seconds   & posterior over the hypotheses of one question         & the question object \eqref{eq:inner_question} \\ \hline
        Days      & priors: where objects are, which doors are shut, what moves & long-term object and topological memory, as Dirichlet counts \\ \hline
        Weeks     & likelihoods: detector and concept models              & false negatives and positives labelled by resolved questions \\ \hline
        Slow      & structure: new categories                             & recurrent, mutually similar ``unknown'' resolutions named by the voice and proposed as a new concept \\ \hline
    \end{tabular}
    \caption{Learning driven by the inner-question loop.}
    \label{tab:inner_learning}
\end{table}

The last row is where the hypothesis of a language for thought is most directly engaged: a cluster of anomalies
that resists every existing concept becomes learnable once it is named. The loop also learns about itself: the
salience weights $\pi_k$ in \eqref{eq:inner_demand} are adjusted by the information each trigger actually yielded,
asking a person is weighed against looking, and loops that repeatedly resolve the same way are compiled into
reflexive rules, so that the number of language-model calls per hour should \emph{decrease} as competence grows.

\subsection{Why a language model: generation versus evaluation}
\label{sec:inner_speech:why_lm}

None of the above requires a person, so the question arises of what a language model contributes to an
architecture that already performs inference and action selection. Our answer is a division of labour. Expected
free energy evaluates policies and Bayesian updating weighs hypotheses, but both operate over a space that has
already been enumerated: each concept agent carries a fixed generative model, policies must be listed before they
can be scored, and Bayesian model reduction can prune structure but not invent it \cite{friston2017curiosity}.
Whatever falls outside the given models remains an unexplained residual. A language model is weak where this
machinery is strong --- its confidences are not calibrated posteriors and it has no access to the sensors --- and
strong where it is weak: it is a cheap, broad generator of hypotheses, explanations and plans. We therefore use it
only as a \emph{proposal distribution}, in a generate-and-test loop whose test is active inference:
\begin{itemize}
    \item \emph{Hypotheses outside the built-in models.} A residual blob can be proposed to be a bag, a toppled
    chair or a robot vacuum, none of which any concept agent represents; the question loop then tests them.
    \item \emph{Amortised policy search.} Instead of enumerating multi-step policies, the model writes a few
    candidate missions --- small programs in the mission language --- which the gate validates and expected free
    energy scores.
    \item \emph{Relational binding across agents.} The agents' representations (signed distance fields, occupancy
    grids, existence log-odds, episodic records) are mutually incompatible; a sentence such as ``the chair that was
    not next to the table in the room behind door~3 yesterday'' binds them in one query no single agent could pose.
    \item \emph{Explanation from few examples.} After a single surprise the model can propose which variable
    mattered (``the door is shut because it is evening''), turning one episode into a testable generalisation.
    \item \emph{Abstraction and naming.} Episodes are compressed into reusable lessons, and recurrent unknowns are
    named, which is what makes them candidates for a new concept (Table~\ref{tab:inner_learning}).
\end{itemize}
Language, rather than an arbitrary learned embedding, matters for three reasons: it is discrete and compositional,
so new hypotheses and plans are assembled from known parts; it is inspectable, so that deterministic gates can parse
and check it and people can audit it; and, once the robot thinks in it, communication with a person comes at no
further cost.

\paragraph{Where the knowledge comes from.}
This raises a difficulty for the motivating thesis. If the inner voice precedes the outer one, the common-sense
knowledge on which the proposer's value rests cannot have arrived through communication --- yet a language model
acquires all of it from human text, that is, from other people's outer voices. The difficulty dissolves once the
\emph{faculty} is distinguished from its \emph{content}. The internalist claim concerns the former: a
computational system for building unbounded hierarchical structures, available for thought before it is used for
talk \cite{hauser2002faculty, berwick2016why}. It makes no claim that what a person knows was acquired without
language. Human knowledge has three sources: core systems for objects, agents, number and space, present before
learning \cite{spelke2007core}; individual experience; and, once outer language exists, cumulative cultural
transmission, which supplies most of what an adult knows \cite{tomasello1999cultural}. Even on the internalist
view, the inner voice of a modern human is filled largely by the outer voices of others.

The architecture has the same three sources, in a different arrangement. The concept agents' hand-built generative
models play the role of \emph{core knowledge}: given before experience and limited to a few domains. The robot's
own counts $n_x(h)$ in \eqref{eq:inner_prior} are \emph{individual experience}. The language model is
\emph{culture}: knowledge distilled from a community's communication and inherited, for this individual, at birth.
The robot thus resembles a child born into a culture more than an early hominin, and its inner voice is Vygotskian
in origin: its format and its content were both learned from outer speech. Two consequences follow. First, the
proposal does not demonstrate the internalist thesis and should not be read as doing so; its claim is functional
--- an inner use of language has epistemic value independently of communication --- and holds whatever the origin of
the language system. Second, the format and the content of the voice can in principle be separated, which turns the
question into an experiment (P7 in Section~\ref{sec:inner_speech:predictions}): a compositional generator over the
robot's own concepts and mission language, trained only on the robot's experience, in the manner of library
learning in program induction \cite{ellis2021dreamcoder}, would supply the format without the culture.

\subsection{Mapping onto the existing architecture}
\label{sec:inner_speech:mapping}

Several components already exist in the language agent. Robot events (mission reports, failures, completions)
already wake the model through a queue, in turns restricted to perceiving, querying memory, stopping and
proposing, never acting. Every action already passes a deterministic gate. The decision to ask or to assume is
already decision-theoretic. Unexplained occupancy is already isolated by a residual agent that acts as the null
hypothesis of the scene, and episodic and long-term object memories are already joined read-only. Missing are the
persistent question objects \eqref{eq:inner_question}, a prediction evaluator implementing
\eqref{eq:inner_update} outside the language model, the demand signal and mode arbitration
\eqref{eq:inner_demand}--\eqref{eq:inner_modes}, the imagined store with its provenance, and the channel through
which bounded spatial priors reach the concept agents.

\subsection{Testable predictions}
\label{sec:inner_speech:predictions}

Consistent with the evaluation principle adopted throughout --- policies compared against each other inside one
closed loop, with body, sensors and environment fixed --- the proposal makes seven falsifiable predictions, each
tested in simulation by moving known objects and introducing unfamiliar ones:
\begin{enumerate}[label=\textbf{P\arabic*}]
    \item The unexplained residual mass per room decreases faster with the focused loop than without it.
    \item The proposer's calibration on resolved questions improves, and $\beta_x$ in \eqref{eq:inner_prior}
    decreases, as the robot's own counts accumulate.
    \item Detector false-negative rates fall after retraining on question-labelled examples.
    \item Inner questions per hour decrease over repeated exposure to the same environment while answer accuracy
    is maintained.
    \item Plans prepared while wandering reduce the latency from surprise to a validated response, and the time to
    notice a deliberately moved object does not grow despite habituation.
    \item With the language-model proposer replaced by a fixed hypothesis set (the existing concept classes and
    ``unknown'') and template-based missions, the loop resolves fewer questions about objects and failures outside
    the built-in models, and acquires new categories more slowly.
    \item A compositional generator trained only on the robot's own experience (format without culture) recovers
    part of the advantage in P6 within a familiar environment, and less of it the further a situation departs from
    what the robot has seen.
\end{enumerate}
A failure of P1 would indicate that the epistemic value of language-mediated questions does not exceed that of
the existing exploration; a failure of P4 would indicate that the loop does not transfer competence to the
automatic layer, which is its principal claim; and a failure of P6 would indicate that the language model adds
nothing a fixed hypothesis set does not, which would remove the case for including it. P7 separates the
contribution of language as a format for thought from that of the knowledge it carries.

\subsection{Limitations}
\label{sec:inner_speech:limitations}

The demand signal \eqref{eq:inner_demand} sums heterogeneous prediction errors whose precisions must themselves be
calibrated, and a poorly scaled term will either suppress the voice or keep it permanently engaged. A single
language model serves the person, the focused loop and the wandering loop; pre-emption keeps the person first but
bounds how much offline reasoning a busy robot can do. The tempered prior \eqref{eq:inner_prior} is a heuristic
combination rather than a derived posterior. The generator of P7 is substantially harder to build than the
ablation of P6, and may have to be approximated by a small model trained on the robot's logs and mission corpus,
whose pre-training would leak some culture back in. Finally, the analogy with human inner speech is functional, not
mechanistic: the proposal is that the same computational role --- serial, compositional self-questioning in the
service of epistemic value --- is useful for a robot, not that the robot's voice resembles the human one in any
further respect.

% and add inner_speech.bib as a second resource:  \addbibresource{inner_speech.bib}
% All keys are new (none collides with the keys main.tex already cites).
% Packages used: amsmath, amssymb, enumitem (already loaded by main.tex) and array, which main.tex does not
% load yet: add \usepackage{array} to its preamble.
% =====================================================================================================

\section{Inner Speech as Epistemic Action}
\label{sec:inner_speech}

The robot described so far uses language in one direction only: a person speaks, and a language agent turns
the request into commands for the controller, grounded in a textual digest of the belief graph. This section
proposes to turn that capability inward. We argue that a language model embedded in an Active Inference
architecture can play the role that inner speech plays in human cognition: a serial, compositional medium in
which the agent poses questions to itself, generates hypotheses, rehearses plans and replays its past, and that
this role follows naturally from the epistemic component of expected free energy. The proposal is stated
together with the constraints that keep such a voice from becoming a source of confabulated belief, and with the
experiments that would falsify it.

\subsection{Motivation: language for thought before language for talk}
\label{sec:inner_speech:motivation}

Two traditions disagree on the origin of inner speech but agree on its function. On the internalist view,
human language evolved first as an internal system for structuring thought, and its use for communication is a
later externalisation \cite{berwick2016why}. On the developmental view, inner speech is social speech that the
child gradually internalises, abbreviates and turns into a tool for self-regulation \cite{vygotsky1986thought}.
Contemporary reviews treat inner speech as functionally involved in planning, problem solving, self-monitoring and
the maintenance of goals in working memory \cite{aldersonday2015inner}, and language more generally as a
cognitive technology that lets a learner manipulate discrete, recombinable tokens rather than raw percepts
\cite{clark1998magic}.

Large language models give a striking instance of the developmental route. They are trained exclusively on
communicative text, yet producing intermediate reasoning in that same medium before answering --- chain-of-thought
--- measurably improves their performance \cite{wei2022chain}. A system built for talk becomes better at thinking
by talking to itself. Our language agent sits at the same point of its own development: it was built to converse
with a person, and the question we ask is what it contributes once its speech is addressed to the robot itself.
Section~\ref{sec:inner_speech:why_lm} returns to what this origin implies for the internalist view.

\subsection{Inner questions as epistemic actions}
\label{sec:inner_speech:questions}

In Active Inference, policies are scored by expected free energy, whose epistemic term rewards the expected
reduction of uncertainty about hidden states \cite{friston2017process, parr2022active}. A question is the
linguistic form of exactly this quantity: it names a hidden variable whose value is uncertain and whose resolution
is expected to matter. We therefore model an inner question as a persistent object
\begin{equation}
    \mathcal{Q} = \big\langle\, e,\; \mathcal{H},\; q(h),\; \mathcal{P},\; b,\; s \,\big\rangle ,
    \label{eq:inner_question}
\end{equation}
where $e$ is the triggering event (for example an obstacle no concept agent explains), $\mathcal{H}$ a finite set
of named hypotheses (``a chair'', ``a box'', ``unknown''), $q(h)$ the current posterior over them, $\mathcal{P}$
the set of pending predictions, $b$ a budget in time and motion, and $s \in \{\text{open},\text{resolved},
\text{abandoned}\}$ its status. Questions form a stack in the sense of a \emph{question under discussion}
\cite{roberts2012information}: ``is it a chair?'' is a sub-question whose answer partially resolves the parent
``what is this?''. The question, not the language model's context window, carries the state of the reasoning from
one step to the next.

\paragraph{One step of the loop.}
Given $\mathcal{Q}$, the language model proposes a test: an epistemic action $a$ (turn towards the anomaly,
approach it, look with the front camera) together with the outcome each hypothesis predicts. The value of the test
is its expected information gain minus its cost,
\begin{equation}
    V(a) \;=\; \underbrace{H\!\left[q(h)\right] \;-\; \mathbb{E}_{p(o \mid a)}\!\left[ H\!\left[q(h \mid o, a)\right]
    \right]}_{I(h;\,o \mid a)} \;-\; c(a),
    \label{eq:inner_value}
\end{equation}
which is the negative of the epistemic part of expected free energy with an explicit action cost. The test is
executed only if $V(a) > 0$ and the action passes the same deterministic gate that already validates every
command the language agent issues. When the observation $o$ arrives, the posterior is updated by the concept
agents' own likelihoods,
\begin{equation}
    q(h \mid o, a) \;\propto\; p(o \mid h, a)\, q(h),
    \label{eq:inner_update}
\end{equation}
and the loop continues, spawns a sub-question, or terminates when
\begin{equation}
    H\!\left[q(h)\right] < \eta \quad\lor\quad \max_a V(a) \le 0 \quad\lor\quad b \text{ exhausted}.
    \label{eq:inner_stop}
\end{equation}
The stopping rule is what separates reasoning from rumination: a loop with no expected information to gain closes,
possibly with the answer ``unknown'', and its trigger is habituated for a refractory period.

\paragraph{The grounding constraint.}
The central design rule is that \emph{the language model proposes and perception disposes}. Text produced by the
model may introduce hypotheses, predictions and tests; it may never move a posterior. Only the likelihood
$p(o \mid h, a)$ in \eqref{eq:inner_update}, computed by the perceptual agents from sensor data, does. This is the
belief-side counterpart of the rule the language agent already enforces for actions, where every proposed command
is validated in C++ against the belief graph before it reaches the controller. Without it, a self-addressed
language loop is a closed system of text conditioned on text and converges on its own confabulations.

\paragraph{The proposer as a calibrated prior.}
The hypotheses the model names, and the probabilities it attaches to them, encode world knowledge the robot has
not acquired by itself (``chairs come with tables''). We treat this as a prior whose weight is earned. Let
$\pi_{\text{LM}}(h \mid x)$ be the proposer's distribution in context $x$ and $n_x(h)$ the robot's own counts of
resolved questions in that context. The prior entering \eqref{eq:inner_update} is
\begin{equation}
    q_0(h \mid x) \;\propto\; \pi_{\text{LM}}(h \mid x)^{\beta_x}\, \big(\alpha + n_x(h)\big),
    \label{eq:inner_prior}
\end{equation}
a tempered product of the proposer and a Dirichlet over experience, where $\beta_x \in [0,1]$ is set from the
proposer's measured calibration in that context (for instance from its Brier score on past questions). Early in
the robot's life the proposer dominates; as counts accumulate, the robot's own experience replaces borrowed
knowledge. The same conjugate form already governs the language agent's decision to ask a person or to assume
that they moved an object, where a Beta posterior over the person's past answers is compared against the costs of
asking and of being wrong.

\subsection{Attention, automaticity and the wandering mind}
\label{sec:inner_speech:modes}

A human driver on an easy road steers without attending to it; the mind is free to plan, imagine or recall, and is
recaptured the instant something unexpected happens \cite{smallwood2015science, christoff2016mind}. In Active
Inference this is a statement about precision: while prediction errors are small and the habitual policy is
reliable, attention --- the precision assigned to sensory prediction errors --- can be withdrawn from the task; a
surprise raises it again \cite{feldman2010attention}. The arbitration between habitual and deliberative control by
their relative uncertainty \cite{daw2005uncertainty} has the same structure.

The architecture already separates the two systems. The controller and the concept agents form an automatic layer
that drives, avoids obstacles and executes missions without the language model; the language agent is the
deliberative voice. We propose to make the allocation of the voice explicit through a scalar \emph{demand} signal,
computed on every digest cycle without the language model,
\begin{equation}
    D_t \;=\; \sum_{k} \pi_k\, \varepsilon_{k,t}^{2} \;+\; \lambda\, \mathbb{1}\!\left[\text{person active}_t\right],
    \label{eq:inner_demand}
\end{equation}
where $\varepsilon_{k,t}$ are normalised prediction errors reported by the agents (unexplained residual occupancy,
flips of existence or door beliefs, localisation uncertainty, stalled or failed mission steps, controller refusals)
and $\pi_k$ their precisions. Three modes follow from two thresholds with asymmetric hysteresis:
\begin{equation}
    m_t =
    \begin{cases}
        \textsc{focused}    & D_t > \theta_{\text{hi}},\\[2pt]
        \textsc{wandering}  & D_\tau < \theta_{\text{lo}} \;\; \forall\, \tau \in [t - T, t],\\[2pt]
        \textsc{monitoring} & \text{otherwise.}
    \end{cases}
    \label{eq:inner_modes}
\end{equation}
Entering \textsc{wandering} requires demand to stay low for a dwell time $T$; leaving it requires a single sample
above $\theta_{\text{hi}}$. The robot is slow to drift off and instant to snap back. Safety never depends on this
arbitration: the controller's reflexes act on their own, and the deliberative voice arrives after the reflex, as
awareness follows braking in the driver.

\paragraph{Focused mode} runs the question loop of Section~\ref{sec:inner_speech:questions} on the triggering
event. \paragraph{Monitoring mode} makes no language-model calls; the voice is primed but silent.

\paragraph{Wandering mode.}
Freed from the task, the voice is seeded with one recent experience chosen by salience: an episode that ended
unresolved, a failure, an object that changed, or an upcoming goal. Each step's output, together with an
associative retrieval from episodic memory, seeds the next, so that the chain drifts as spontaneous thought does
\cite{christoff2016mind}. The voice has access only to memory queries and to the architecture's imagination tools,
which simulate observations and outcomes on a copy of the belief graph; it can neither act nor write beliefs. Its
products are of three kinds: counterfactuals over past episodes (``had the door been open, the crossing would have
taken 40\,s less''), hypotheses to be tested later, and prospective plans (``if the door is shut on the way back,
the route through the next room is shorter''). Functionally this is the default network's role in simulating
past and future \cite{buckner2008brain}, and the offline replay and model reduction that Active Inference assigns
to rest and sleep \cite{friston2017curiosity, hobson2012waking}. Replay is prioritised by the expected gain of the
memory it revisits \cite{mattar2018prioritized}.

\paragraph{Source monitoring.}
Humans occasionally misremember imagined events as perceived ones \cite{johnson1993source}. The robot must not.
Everything produced while wandering is stored with an \emph{imagined} provenance, separate from the factual
episodic record, and enters the learning loop only as a hypothesis to be tested, never as evidence.

\subsection{What the return to reality gives the automatic layer}
\label{sec:inner_speech:return}

The transition from \textsc{wandering} to \textsc{focused} cancels the chain in flight (its state is kept, so that
it may resume if demand subsides) and gives the full voice to the surprise. It serves the automatic layer at three
timescales.
\begin{enumerate}
    \item \textbf{Decision.} A surprise means the automatic layer has left its domain of competence, where its only
    responses are to retry, stop or fail. The voice adds re-interpretation of the scene, arbitration between goals,
    re-planning through remembered topology, and recourse to a person.
    \item \textbf{Attention.} The voice can raise precision where it matters: orienting sensors, increasing a
    detector's rate in a region, or passing a hypothesis to the concept agents as a spatial prior. Such a prior has
    bounded weight; it may accelerate the recognition of an object that is present but may not create one that is
    not, preserving the grounding constraint.
    \item \textbf{Teaching.} Each resolved surprise is a labelled example the automatic layer could not have
    produced alone: a residual later identified as a chair is a false negative of the chair detector; a door that
    is repeatedly shut in the evening becomes a prior; a recovery that repeatedly succeeds becomes a branch of the
    controller's behaviour tree. The domain of competence grows, and the situations that require the voice become
    rarer, which is the classical progression from effortful to autonomous skill \cite{fitts1967human}.
\end{enumerate}
Plans prepared while wandering act as a cache: when the anticipated situation occurs, the focused voice verifies a
plan instead of composing one. The return ends with an explicit release --- spatial priors cleared, safety margins
restored, the automatic layer confirmed back within its competence --- without which every surprise would leave
the robot permanently vigilant.

\subsection{Relation to chain-of-thought reasoning}
\label{sec:inner_speech:cot}

Chain-of-thought is inner speech in the functional sense, but of a specific kind, and the comparison sharpens the
design (Table~\ref{tab:inner_cot}). Plain chain-of-thought is convergent and ungrounded: tokens condition on
tokens until an answer is produced, and the trace is discarded. Interleaving reasoning with action and observation
grounds it \cite{yao2023react}; our focused loop is of this kind, with the additional property that the
observation, through \eqref{eq:inner_update}, rather than the next token, determines where the reasoning goes.
The wandering voice is closer to plain chain-of-thought in its isolation from perception, but divergent rather
than convergent, interruptible, and persistent in its products.

\begin{table}[ht]
    \centering
    \small
    \begin{tabular}{|>{\raggedright\arraybackslash}p{0.19\linewidth}|>{\raggedright\arraybackslash}p{0.22\linewidth}|>{\raggedright\arraybackslash}p{0.24\linewidth}|>{\raggedright\arraybackslash}p{0.24\linewidth}|}
        \hline
                                & \textbf{Chain-of-thought} & \textbf{Focused voice} & \textbf{Wandering voice} \\ \hline
        Goal                    & answer a given question   & resolve a given surprise & none given; seed by salience \\ \hline
        Dynamics                & convergent                & convergent             & divergent, associative \\ \hline
        Contact with the world  & none                      & thought--action--observation & none (imagination only) \\ \hline
        Persistence             & discarded                 & question objects, beliefs, lessons & imagined store, hypotheses \\ \hline
        Interruptible           & no                        & yes                    & yes, resumable \\ \hline
        Latency budget          & n/a                       & $\sim$1\,s, no extended thinking & unbounded, extended thinking \\ \hline
    \end{tabular}
    \caption{Chain-of-thought compared with the two modes of the proposed inner voice.}
    \label{tab:inner_cot}
\end{table}

Two practical consequences follow. First, \emph{think when there is time, not when there is pressure}. In the
deployed language agent extended thinking is disabled, because a thinking-mode model spent 3\,000--6\,000 tokens
(40--80\,s) before each tool call; that is incompatible with the focused mode but costs nothing in the wandering
mode, where it can be enabled with a large budget and pre-empted at any moment. Slow reasoning is done in advance
and fast reasoning uses its results. Second, chain-of-thought traces are known not always to reflect the
computation that produced the answer \cite{turpin2023language}, as human verbal reports do not always reflect the
causes of behaviour \cite{nisbett1977telling}. The faithful record of the robot's reasoning is therefore the
structured one --- hypotheses, predictions, observations and posterior updates --- not the text. Finally, wandering
chains whose hypotheses were later confirmed by perception form a corpus of reasoning verified against the world,
which could be used to fine-tune the local model in the manner of self-taught reasoning \cite{zelikman2022star},
with perception, rather than a known answer, as the filter.

\subsection{Consequences for belief and learning}
\label{sec:inner_speech:learning}

The loop changes how beliefs evolve in two ways. Within an episode, uncertainty is reduced where surprise
occurred rather than where the sensors happened to point, which is efficient but leaves the rest of the map less
examined and makes the robot's own data a biased sample; a decaying novelty term \cite{schwartenbeck2019curiosity}
and importance weighting of learned statistics are required to correct this. Across episodes, every closed
question is a labelled example, and learning proceeds at four timescales (Table~\ref{tab:inner_learning}).

\begin{table}[ht]
    \centering
    \small
    \begin{tabular}{|>{\raggedright\arraybackslash}p{0.14\linewidth}|>{\raggedright\arraybackslash}p{0.37\linewidth}|>{\raggedright\arraybackslash}p{0.37\linewidth}|}
        \hline
        \textbf{Timescale} & \textbf{What changes} & \textbf{Where, in this architecture} \\ \hline
        Seconds   & posterior over the hypotheses of one question         & the question object \eqref{eq:inner_question} \\ \hline
        Days      & priors: where objects are, which doors are shut, what moves & long-term object and topological memory, as Dirichlet counts \\ \hline
        Weeks     & likelihoods: detector and concept models              & false negatives and positives labelled by resolved questions \\ \hline
        Slow      & structure: new categories                             & recurrent, mutually similar ``unknown'' resolutions named by the voice and proposed as a new concept \\ \hline
    \end{tabular}
    \caption{Learning driven by the inner-question loop.}
    \label{tab:inner_learning}
\end{table}

The last row is where the hypothesis of a language for thought is most directly engaged: a cluster of anomalies
that resists every existing concept becomes learnable once it is named. The loop also learns about itself: the
salience weights $\pi_k$ in \eqref{eq:inner_demand} are adjusted by the information each trigger actually yielded,
asking a person is weighed against looking, and loops that repeatedly resolve the same way are compiled into
reflexive rules, so that the number of language-model calls per hour should \emph{decrease} as competence grows.

\subsection{Why a language model: generation versus evaluation}
\label{sec:inner_speech:why_lm}

None of the above requires a person, so the question arises of what a language model contributes to an
architecture that already performs inference and action selection. Our answer is a division of labour. Expected
free energy evaluates policies and Bayesian updating weighs hypotheses, but both operate over a space that has
already been enumerated: each concept agent carries a fixed generative model, policies must be listed before they
can be scored, and Bayesian model reduction can prune structure but not invent it \cite{friston2017curiosity}.
Whatever falls outside the given models remains an unexplained residual. A language model is weak where this
machinery is strong --- its confidences are not calibrated posteriors and it has no access to the sensors --- and
strong where it is weak: it is a cheap, broad generator of hypotheses, explanations and plans. We therefore use it
only as a \emph{proposal distribution}, in a generate-and-test loop whose test is active inference:
\begin{itemize}
    \item \emph{Hypotheses outside the built-in models.} A residual blob can be proposed to be a bag, a toppled
    chair or a robot vacuum, none of which any concept agent represents; the question loop then tests them.
    \item \emph{Amortised policy search.} Instead of enumerating multi-step policies, the model writes a few
    candidate missions --- small programs in the mission language --- which the gate validates and expected free
    energy scores.
    \item \emph{Relational binding across agents.} The agents' representations (signed distance fields, occupancy
    grids, existence log-odds, episodic records) are mutually incompatible; a sentence such as ``the chair that was
    not next to the table in the room behind door~3 yesterday'' binds them in one query no single agent could pose.
    \item \emph{Explanation from few examples.} After a single surprise the model can propose which variable
    mattered (``the door is shut because it is evening''), turning one episode into a testable generalisation.
    \item \emph{Abstraction and naming.} Episodes are compressed into reusable lessons, and recurrent unknowns are
    named, which is what makes them candidates for a new concept (Table~\ref{tab:inner_learning}).
\end{itemize}
Language, rather than an arbitrary learned embedding, matters for three reasons: it is discrete and compositional,
so new hypotheses and plans are assembled from known parts; it is inspectable, so that deterministic gates can parse
and check it and people can audit it; and, once the robot thinks in it, communication with a person comes at no
further cost.

\paragraph{Where the knowledge comes from.}
This raises a difficulty for the motivating thesis. If the inner voice precedes the outer one, the common-sense
knowledge on which the proposer's value rests cannot have arrived through communication --- yet a language model
acquires all of it from human text, that is, from other people's outer voices. The difficulty dissolves once the
\emph{faculty} is distinguished from its \emph{content}. The internalist claim concerns the former: a
computational system for building unbounded hierarchical structures, available for thought before it is used for
talk \cite{hauser2002faculty, berwick2016why}. It makes no claim that what a person knows was acquired without
language. Human knowledge has three sources: core systems for objects, agents, number and space, present before
learning \cite{spelke2007core}; individual experience; and, once outer language exists, cumulative cultural
transmission, which supplies most of what an adult knows \cite{tomasello1999cultural}. Even on the internalist
view, the inner voice of a modern human is filled largely by the outer voices of others.

The architecture has the same three sources, in a different arrangement. The concept agents' hand-built generative
models play the role of \emph{core knowledge}: given before experience and limited to a few domains. The robot's
own counts $n_x(h)$ in \eqref{eq:inner_prior} are \emph{individual experience}. The language model is
\emph{culture}: knowledge distilled from a community's communication and inherited, for this individual, at birth.
The robot thus resembles a child born into a culture more than an early hominin, and its inner voice is Vygotskian
in origin: its format and its content were both learned from outer speech. Two consequences follow. First, the
proposal does not demonstrate the internalist thesis and should not be read as doing so; its claim is functional
--- an inner use of language has epistemic value independently of communication --- and holds whatever the origin of
the language system. Second, the format and the content of the voice can in principle be separated, which turns the
question into an experiment (P7 in Section~\ref{sec:inner_speech:predictions}): a compositional generator over the
robot's own concepts and mission language, trained only on the robot's experience, in the manner of library
learning in program induction \cite{ellis2021dreamcoder}, would supply the format without the culture.

\subsection{Mapping onto the existing architecture}
\label{sec:inner_speech:mapping}

Several components already exist in the language agent. Robot events (mission reports, failures, completions)
already wake the model through a queue, in turns restricted to perceiving, querying memory, stopping and
proposing, never acting. Every action already passes a deterministic gate. The decision to ask or to assume is
already decision-theoretic. Unexplained occupancy is already isolated by a residual agent that acts as the null
hypothesis of the scene, and episodic and long-term object memories are already joined read-only. Missing are the
persistent question objects \eqref{eq:inner_question}, a prediction evaluator implementing
\eqref{eq:inner_update} outside the language model, the demand signal and mode arbitration
\eqref{eq:inner_demand}--\eqref{eq:inner_modes}, the imagined store with its provenance, and the channel through
which bounded spatial priors reach the concept agents.

\subsection{Testable predictions}
\label{sec:inner_speech:predictions}

Consistent with the evaluation principle adopted throughout --- policies compared against each other inside one
closed loop, with body, sensors and environment fixed --- the proposal makes seven falsifiable predictions, each
tested in simulation by moving known objects and introducing unfamiliar ones:
\begin{enumerate}[label=\textbf{P\arabic*}]
    \item The unexplained residual mass per room decreases faster with the focused loop than without it.
    \item The proposer's calibration on resolved questions improves, and $\beta_x$ in \eqref{eq:inner_prior}
    decreases, as the robot's own counts accumulate.
    \item Detector false-negative rates fall after retraining on question-labelled examples.
    \item Inner questions per hour decrease over repeated exposure to the same environment while answer accuracy
    is maintained.
    \item Plans prepared while wandering reduce the latency from surprise to a validated response, and the time to
    notice a deliberately moved object does not grow despite habituation.
    \item With the language-model proposer replaced by a fixed hypothesis set (the existing concept classes and
    ``unknown'') and template-based missions, the loop resolves fewer questions about objects and failures outside
    the built-in models, and acquires new categories more slowly.
    \item A compositional generator trained only on the robot's own experience (format without culture) recovers
    part of the advantage in P6 within a familiar environment, and less of it the further a situation departs from
    what the robot has seen.
\end{enumerate}
A failure of P1 would indicate that the epistemic value of language-mediated questions does not exceed that of
the existing exploration; a failure of P4 would indicate that the loop does not transfer competence to the
automatic layer, which is its principal claim; and a failure of P6 would indicate that the language model adds
nothing a fixed hypothesis set does not, which would remove the case for including it. P7 separates the
contribution of language as a format for thought from that of the knowledge it carries.

\subsection{Limitations}
\label{sec:inner_speech:limitations}

The demand signal \eqref{eq:inner_demand} sums heterogeneous prediction errors whose precisions must themselves be
calibrated, and a poorly scaled term will either suppress the voice or keep it permanently engaged. A single
language model serves the person, the focused loop and the wandering loop; pre-emption keeps the person first but
bounds how much offline reasoning a busy robot can do. The tempered prior \eqref{eq:inner_prior} is a heuristic
combination rather than a derived posterior. The generator of P7 is substantially harder to build than the
ablation of P6, and may have to be approximated by a small model trained on the robot's logs and mission corpus,
whose pre-training would leak some culture back in. Finally, the analogy with human inner speech is functional, not
mechanistic: the proposal is that the same computational role --- serial, compositional self-questioning in the
service of epistemic value --- is useful for a robot, not that the robot's voice resembles the human one in any
further respect.

% and add inner_speech.bib as a second resource:  \addbibresource{inner_speech.bib}
% All keys are new (none collides with the keys main.tex already cites).
% Packages used: amsmath, amssymb, enumitem (already loaded by main.tex) and array, which main.tex does not
% load yet: add \usepackage{array} to its preamble.
% =====================================================================================================

\section{Inner Speech as Epistemic Action}
\label{sec:inner_speech}

The robot described so far uses language in one direction only: a person speaks, and a language agent turns
the request into commands for the controller, grounded in a textual digest of the belief graph. This section
proposes to turn that capability inward. We argue that a language model embedded in an Active Inference
architecture can play the role that inner speech plays in human cognition: a serial, compositional medium in
which the agent poses questions to itself, generates hypotheses, rehearses plans and replays its past, and that
this role follows naturally from the epistemic component of expected free energy. The proposal is stated
together with the constraints that keep such a voice from becoming a source of confabulated belief, and with the
experiments that would falsify it.

\subsection{Motivation: language for thought before language for talk}
\label{sec:inner_speech:motivation}

Two traditions disagree on the origin of inner speech but agree on its function. On the internalist view,
human language evolved first as an internal system for structuring thought, and its use for communication is a
later externalisation \cite{berwick2016why}. On the developmental view, inner speech is social speech that the
child gradually internalises, abbreviates and turns into a tool for self-regulation \cite{vygotsky1986thought}.
Contemporary reviews treat inner speech as functionally involved in planning, problem solving, self-monitoring and
the maintenance of goals in working memory \cite{aldersonday2015inner}, and language more generally as a
cognitive technology that lets a learner manipulate discrete, recombinable tokens rather than raw percepts
\cite{clark1998magic}.

Large language models give a striking instance of the developmental route. They are trained exclusively on
communicative text, yet producing intermediate reasoning in that same medium before answering --- chain-of-thought
--- measurably improves their performance \cite{wei2022chain}. A system built for talk becomes better at thinking
by talking to itself. Our language agent sits at the same point of its own development: it was built to converse
with a person, and the question we ask is what it contributes once its speech is addressed to the robot itself.
Section~\ref{sec:inner_speech:why_lm} returns to what this origin implies for the internalist view.

\subsection{Inner questions as epistemic actions}
\label{sec:inner_speech:questions}

In Active Inference, policies are scored by expected free energy, whose epistemic term rewards the expected
reduction of uncertainty about hidden states \cite{friston2017process, parr2022active}. A question is the
linguistic form of exactly this quantity: it names a hidden variable whose value is uncertain and whose resolution
is expected to matter. We therefore model an inner question as a persistent object
\begin{equation}
    \mathcal{Q} = \big\langle\, e,\; \mathcal{H},\; q(h),\; \mathcal{P},\; b,\; s \,\big\rangle ,
    \label{eq:inner_question}
\end{equation}
where $e$ is the triggering event (for example an obstacle no concept agent explains), $\mathcal{H}$ a finite set
of named hypotheses (``a chair'', ``a box'', ``unknown''), $q(h)$ the current posterior over them, $\mathcal{P}$
the set of pending predictions, $b$ a budget in time and motion, and $s \in \{\text{open},\text{resolved},
\text{abandoned}\}$ its status. Questions form a stack in the sense of a \emph{question under discussion}
\cite{roberts2012information}: ``is it a chair?'' is a sub-question whose answer partially resolves the parent
``what is this?''. The question, not the language model's context window, carries the state of the reasoning from
one step to the next.

\paragraph{One step of the loop.}
Given $\mathcal{Q}$, the language model proposes a test: an epistemic action $a$ (turn towards the anomaly,
approach it, look with the front camera) together with the outcome each hypothesis predicts. The value of the test
is its expected information gain minus its cost,
\begin{equation}
    V(a) \;=\; \underbrace{H\!\left[q(h)\right] \;-\; \mathbb{E}_{p(o \mid a)}\!\left[ H\!\left[q(h \mid o, a)\right]
    \right]}_{I(h;\,o \mid a)} \;-\; c(a),
    \label{eq:inner_value}
\end{equation}
which is the negative of the epistemic part of expected free energy with an explicit action cost. The test is
executed only if $V(a) > 0$ and the action passes the same deterministic gate that already validates every
command the language agent issues. When the observation $o$ arrives, the posterior is updated by the concept
agents' own likelihoods,
\begin{equation}
    q(h \mid o, a) \;\propto\; p(o \mid h, a)\, q(h),
    \label{eq:inner_update}
\end{equation}
and the loop continues, spawns a sub-question, or terminates when
\begin{equation}
    H\!\left[q(h)\right] < \eta \quad\lor\quad \max_a V(a) \le 0 \quad\lor\quad b \text{ exhausted}.
    \label{eq:inner_stop}
\end{equation}
The stopping rule is what separates reasoning from rumination: a loop with no expected information to gain closes,
possibly with the answer ``unknown'', and its trigger is habituated for a refractory period.

\paragraph{The grounding constraint.}
The central design rule is that \emph{the language model proposes and perception disposes}. Text produced by the
model may introduce hypotheses, predictions and tests; it may never move a posterior. Only the likelihood
$p(o \mid h, a)$ in \eqref{eq:inner_update}, computed by the perceptual agents from sensor data, does. This is the
belief-side counterpart of the rule the language agent already enforces for actions, where every proposed command
is validated in C++ against the belief graph before it reaches the controller. Without it, a self-addressed
language loop is a closed system of text conditioned on text and converges on its own confabulations.

\paragraph{The proposer as a calibrated prior.}
The hypotheses the model names, and the probabilities it attaches to them, encode world knowledge the robot has
not acquired by itself (``chairs come with tables''). We treat this as a prior whose weight is earned. Let
$\pi_{\text{LM}}(h \mid x)$ be the proposer's distribution in context $x$ and $n_x(h)$ the robot's own counts of
resolved questions in that context. The prior entering \eqref{eq:inner_update} is
\begin{equation}
    q_0(h \mid x) \;\propto\; \pi_{\text{LM}}(h \mid x)^{\beta_x}\, \big(\alpha + n_x(h)\big),
    \label{eq:inner_prior}
\end{equation}
a tempered product of the proposer and a Dirichlet over experience, where $\beta_x \in [0,1]$ is set from the
proposer's measured calibration in that context (for instance from its Brier score on past questions). Early in
the robot's life the proposer dominates; as counts accumulate, the robot's own experience replaces borrowed
knowledge. The same conjugate form already governs the language agent's decision to ask a person or to assume
that they moved an object, where a Beta posterior over the person's past answers is compared against the costs of
asking and of being wrong.

\subsection{Attention, automaticity and the wandering mind}
\label{sec:inner_speech:modes}

A human driver on an easy road steers without attending to it; the mind is free to plan, imagine or recall, and is
recaptured the instant something unexpected happens \cite{smallwood2015science, christoff2016mind}. In Active
Inference this is a statement about precision: while prediction errors are small and the habitual policy is
reliable, attention --- the precision assigned to sensory prediction errors --- can be withdrawn from the task; a
surprise raises it again \cite{feldman2010attention}. The arbitration between habitual and deliberative control by
their relative uncertainty \cite{daw2005uncertainty} has the same structure.

The architecture already separates the two systems. The controller and the concept agents form an automatic layer
that drives, avoids obstacles and executes missions without the language model; the language agent is the
deliberative voice. We propose to make the allocation of the voice explicit through a scalar \emph{demand} signal,
computed on every digest cycle without the language model,
\begin{equation}
    D_t \;=\; \sum_{k} \pi_k\, \varepsilon_{k,t}^{2} \;+\; \lambda\, \mathbb{1}\!\left[\text{person active}_t\right],
    \label{eq:inner_demand}
\end{equation}
where $\varepsilon_{k,t}$ are normalised prediction errors reported by the agents (unexplained residual occupancy,
flips of existence or door beliefs, localisation uncertainty, stalled or failed mission steps, controller refusals)
and $\pi_k$ their precisions. Three modes follow from two thresholds with asymmetric hysteresis:
\begin{equation}
    m_t =
    \begin{cases}
        \textsc{focused}    & D_t > \theta_{\text{hi}},\\[2pt]
        \textsc{wandering}  & D_\tau < \theta_{\text{lo}} \;\; \forall\, \tau \in [t - T, t],\\[2pt]
        \textsc{monitoring} & \text{otherwise.}
    \end{cases}
    \label{eq:inner_modes}
\end{equation}
Entering \textsc{wandering} requires demand to stay low for a dwell time $T$; leaving it requires a single sample
above $\theta_{\text{hi}}$. The robot is slow to drift off and instant to snap back. Safety never depends on this
arbitration: the controller's reflexes act on their own, and the deliberative voice arrives after the reflex, as
awareness follows braking in the driver.

\paragraph{Focused mode} runs the question loop of Section~\ref{sec:inner_speech:questions} on the triggering
event. \paragraph{Monitoring mode} makes no language-model calls; the voice is primed but silent.

\paragraph{Wandering mode.}
Freed from the task, the voice is seeded with one recent experience chosen by salience: an episode that ended
unresolved, a failure, an object that changed, or an upcoming goal. Each step's output, together with an
associative retrieval from episodic memory, seeds the next, so that the chain drifts as spontaneous thought does
\cite{christoff2016mind}. The voice has access only to memory queries and to the architecture's imagination tools,
which simulate observations and outcomes on a copy of the belief graph; it can neither act nor write beliefs. Its
products are of three kinds: counterfactuals over past episodes (``had the door been open, the crossing would have
taken 40\,s less''), hypotheses to be tested later, and prospective plans (``if the door is shut on the way back,
the route through the next room is shorter''). Functionally this is the default network's role in simulating
past and future \cite{buckner2008brain}, and the offline replay and model reduction that Active Inference assigns
to rest and sleep \cite{friston2017curiosity, hobson2012waking}. Replay is prioritised by the expected gain of the
memory it revisits \cite{mattar2018prioritized}.

\paragraph{Source monitoring.}
Humans occasionally misremember imagined events as perceived ones \cite{johnson1993source}. The robot must not.
Everything produced while wandering is stored with an \emph{imagined} provenance, separate from the factual
episodic record, and enters the learning loop only as a hypothesis to be tested, never as evidence.

\subsection{What the return to reality gives the automatic layer}
\label{sec:inner_speech:return}

The transition from \textsc{wandering} to \textsc{focused} cancels the chain in flight (its state is kept, so that
it may resume if demand subsides) and gives the full voice to the surprise. It serves the automatic layer at three
timescales.
\begin{enumerate}
    \item \textbf{Decision.} A surprise means the automatic layer has left its domain of competence, where its only
    responses are to retry, stop or fail. The voice adds re-interpretation of the scene, arbitration between goals,
    re-planning through remembered topology, and recourse to a person.
    \item \textbf{Attention.} The voice can raise precision where it matters: orienting sensors, increasing a
    detector's rate in a region, or passing a hypothesis to the concept agents as a spatial prior. Such a prior has
    bounded weight; it may accelerate the recognition of an object that is present but may not create one that is
    not, preserving the grounding constraint.
    \item \textbf{Teaching.} Each resolved surprise is a labelled example the automatic layer could not have
    produced alone: a residual later identified as a chair is a false negative of the chair detector; a door that
    is repeatedly shut in the evening becomes a prior; a recovery that repeatedly succeeds becomes a branch of the
    controller's behaviour tree. The domain of competence grows, and the situations that require the voice become
    rarer, which is the classical progression from effortful to autonomous skill \cite{fitts1967human}.
\end{enumerate}
Plans prepared while wandering act as a cache: when the anticipated situation occurs, the focused voice verifies a
plan instead of composing one. The return ends with an explicit release --- spatial priors cleared, safety margins
restored, the automatic layer confirmed back within its competence --- without which every surprise would leave
the robot permanently vigilant.

\subsection{Relation to chain-of-thought reasoning}
\label{sec:inner_speech:cot}

Chain-of-thought is inner speech in the functional sense, but of a specific kind, and the comparison sharpens the
design (Table~\ref{tab:inner_cot}). Plain chain-of-thought is convergent and ungrounded: tokens condition on
tokens until an answer is produced, and the trace is discarded. Interleaving reasoning with action and observation
grounds it \cite{yao2023react}; our focused loop is of this kind, with the additional property that the
observation, through \eqref{eq:inner_update}, rather than the next token, determines where the reasoning goes.
The wandering voice is closer to plain chain-of-thought in its isolation from perception, but divergent rather
than convergent, interruptible, and persistent in its products.

\begin{table}[ht]
    \centering
    \small
    \begin{tabular}{|>{\raggedright\arraybackslash}p{0.19\linewidth}|>{\raggedright\arraybackslash}p{0.22\linewidth}|>{\raggedright\arraybackslash}p{0.24\linewidth}|>{\raggedright\arraybackslash}p{0.24\linewidth}|}
        \hline
                                & \textbf{Chain-of-thought} & \textbf{Focused voice} & \textbf{Wandering voice} \\ \hline
        Goal                    & answer a given question   & resolve a given surprise & none given; seed by salience \\ \hline
        Dynamics                & convergent                & convergent             & divergent, associative \\ \hline
        Contact with the world  & none                      & thought--action--observation & none (imagination only) \\ \hline
        Persistence             & discarded                 & question objects, beliefs, lessons & imagined store, hypotheses \\ \hline
        Interruptible           & no                        & yes                    & yes, resumable \\ \hline
        Latency budget          & n/a                       & $\sim$1\,s, no extended thinking & unbounded, extended thinking \\ \hline
    \end{tabular}
    \caption{Chain-of-thought compared with the two modes of the proposed inner voice.}
    \label{tab:inner_cot}
\end{table}

Two practical consequences follow. First, \emph{think when there is time, not when there is pressure}. In the
deployed language agent extended thinking is disabled, because a thinking-mode model spent 3\,000--6\,000 tokens
(40--80\,s) before each tool call; that is incompatible with the focused mode but costs nothing in the wandering
mode, where it can be enabled with a large budget and pre-empted at any moment. Slow reasoning is done in advance
and fast reasoning uses its results. Second, chain-of-thought traces are known not always to reflect the
computation that produced the answer \cite{turpin2023language}, as human verbal reports do not always reflect the
causes of behaviour \cite{nisbett1977telling}. The faithful record of the robot's reasoning is therefore the
structured one --- hypotheses, predictions, observations and posterior updates --- not the text. Finally, wandering
chains whose hypotheses were later confirmed by perception form a corpus of reasoning verified against the world,
which could be used to fine-tune the local model in the manner of self-taught reasoning \cite{zelikman2022star},
with perception, rather than a known answer, as the filter.

\subsection{Consequences for belief and learning}
\label{sec:inner_speech:learning}

The loop changes how beliefs evolve in two ways. Within an episode, uncertainty is reduced where surprise
occurred rather than where the sensors happened to point, which is efficient but leaves the rest of the map less
examined and makes the robot's own data a biased sample; a decaying novelty term \cite{schwartenbeck2019curiosity}
and importance weighting of learned statistics are required to correct this. Across episodes, every closed
question is a labelled example, and learning proceeds at four timescales (Table~\ref{tab:inner_learning}).

\begin{table}[ht]
    \centering
    \small
    \begin{tabular}{|>{\raggedright\arraybackslash}p{0.14\linewidth}|>{\raggedright\arraybackslash}p{0.37\linewidth}|>{\raggedright\arraybackslash}p{0.37\linewidth}|}
        \hline
        \textbf{Timescale} & \textbf{What changes} & \textbf{Where, in this architecture} \\ \hline
        Seconds   & posterior over the hypotheses of one question         & the question object \eqref{eq:inner_question} \\ \hline
        Days      & priors: where objects are, which doors are shut, what moves & long-term object and topological memory, as Dirichlet counts \\ \hline
        Weeks     & likelihoods: detector and concept models              & false negatives and positives labelled by resolved questions \\ \hline
        Slow      & structure: new categories                             & recurrent, mutually similar ``unknown'' resolutions named by the voice and proposed as a new concept \\ \hline
    \end{tabular}
    \caption{Learning driven by the inner-question loop.}
    \label{tab:inner_learning}
\end{table}

The last row is where the hypothesis of a language for thought is most directly engaged: a cluster of anomalies
that resists every existing concept becomes learnable once it is named. The loop also learns about itself: the
salience weights $\pi_k$ in \eqref{eq:inner_demand} are adjusted by the information each trigger actually yielded,
asking a person is weighed against looking, and loops that repeatedly resolve the same way are compiled into
reflexive rules, so that the number of language-model calls per hour should \emph{decrease} as competence grows.

\subsection{Why a language model: generation versus evaluation}
\label{sec:inner_speech:why_lm}

None of the above requires a person, so the question arises of what a language model contributes to an
architecture that already performs inference and action selection. Our answer is a division of labour. Expected
free energy evaluates policies and Bayesian updating weighs hypotheses, but both operate over a space that has
already been enumerated: each concept agent carries a fixed generative model, policies must be listed before they
can be scored, and Bayesian model reduction can prune structure but not invent it \cite{friston2017curiosity}.
Whatever falls outside the given models remains an unexplained residual. A language model is weak where this
machinery is strong --- its confidences are not calibrated posteriors and it has no access to the sensors --- and
strong where it is weak: it is a cheap, broad generator of hypotheses, explanations and plans. We therefore use it
only as a \emph{proposal distribution}, in a generate-and-test loop whose test is active inference:
\begin{itemize}
    \item \emph{Hypotheses outside the built-in models.} A residual blob can be proposed to be a bag, a toppled
    chair or a robot vacuum, none of which any concept agent represents; the question loop then tests them.
    \item \emph{Amortised policy search.} Instead of enumerating multi-step policies, the model writes a few
    candidate missions --- small programs in the mission language --- which the gate validates and expected free
    energy scores.
    \item \emph{Relational binding across agents.} The agents' representations (signed distance fields, occupancy
    grids, existence log-odds, episodic records) are mutually incompatible; a sentence such as ``the chair that was
    not next to the table in the room behind door~3 yesterday'' binds them in one query no single agent could pose.
    \item \emph{Explanation from few examples.} After a single surprise the model can propose which variable
    mattered (``the door is shut because it is evening''), turning one episode into a testable generalisation.
    \item \emph{Abstraction and naming.} Episodes are compressed into reusable lessons, and recurrent unknowns are
    named, which is what makes them candidates for a new concept (Table~\ref{tab:inner_learning}).
\end{itemize}
Language, rather than an arbitrary learned embedding, matters for three reasons: it is discrete and compositional,
so new hypotheses and plans are assembled from known parts; it is inspectable, so that deterministic gates can parse
and check it and people can audit it; and, once the robot thinks in it, communication with a person comes at no
further cost.

\paragraph{Where the knowledge comes from.}
This raises a difficulty for the motivating thesis. If the inner voice precedes the outer one, the common-sense
knowledge on which the proposer's value rests cannot have arrived through communication --- yet a language model
acquires all of it from human text, that is, from other people's outer voices. The difficulty dissolves once the
\emph{faculty} is distinguished from its \emph{content}. The internalist claim concerns the former: a
computational system for building unbounded hierarchical structures, available for thought before it is used for
talk \cite{hauser2002faculty, berwick2016why}. It makes no claim that what a person knows was acquired without
language. Human knowledge has three sources: core systems for objects, agents, number and space, present before
learning \cite{spelke2007core}; individual experience; and, once outer language exists, cumulative cultural
transmission, which supplies most of what an adult knows \cite{tomasello1999cultural}. Even on the internalist
view, the inner voice of a modern human is filled largely by the outer voices of others.

The architecture has the same three sources, in a different arrangement. The concept agents' hand-built generative
models play the role of \emph{core knowledge}: given before experience and limited to a few domains. The robot's
own counts $n_x(h)$ in \eqref{eq:inner_prior} are \emph{individual experience}. The language model is
\emph{culture}: knowledge distilled from a community's communication and inherited, for this individual, at birth.
The robot thus resembles a child born into a culture more than an early hominin, and its inner voice is Vygotskian
in origin: its format and its content were both learned from outer speech. Two consequences follow. First, the
proposal does not demonstrate the internalist thesis and should not be read as doing so; its claim is functional
--- an inner use of language has epistemic value independently of communication --- and holds whatever the origin of
the language system. Second, the format and the content of the voice can in principle be separated, which turns the
question into an experiment (P7 in Section~\ref{sec:inner_speech:predictions}): a compositional generator over the
robot's own concepts and mission language, trained only on the robot's experience, in the manner of library
learning in program induction \cite{ellis2021dreamcoder}, would supply the format without the culture.

\subsection{Mapping onto the existing architecture}
\label{sec:inner_speech:mapping}

Several components already exist in the language agent. Robot events (mission reports, failures, completions)
already wake the model through a queue, in turns restricted to perceiving, querying memory, stopping and
proposing, never acting. Every action already passes a deterministic gate. The decision to ask or to assume is
already decision-theoretic. Unexplained occupancy is already isolated by a residual agent that acts as the null
hypothesis of the scene, and episodic and long-term object memories are already joined read-only. Missing are the
persistent question objects \eqref{eq:inner_question}, a prediction evaluator implementing
\eqref{eq:inner_update} outside the language model, the demand signal and mode arbitration
\eqref{eq:inner_demand}--\eqref{eq:inner_modes}, the imagined store with its provenance, and the channel through
which bounded spatial priors reach the concept agents.

\subsection{Testable predictions}
\label{sec:inner_speech:predictions}

Consistent with the evaluation principle adopted throughout --- policies compared against each other inside one
closed loop, with body, sensors and environment fixed --- the proposal makes seven falsifiable predictions, each
tested in simulation by moving known objects and introducing unfamiliar ones:
\begin{enumerate}[label=\textbf{P\arabic*}]
    \item The unexplained residual mass per room decreases faster with the focused loop than without it.
    \item The proposer's calibration on resolved questions improves, and $\beta_x$ in \eqref{eq:inner_prior}
    decreases, as the robot's own counts accumulate.
    \item Detector false-negative rates fall after retraining on question-labelled examples.
    \item Inner questions per hour decrease over repeated exposure to the same environment while answer accuracy
    is maintained.
    \item Plans prepared while wandering reduce the latency from surprise to a validated response, and the time to
    notice a deliberately moved object does not grow despite habituation.
    \item With the language-model proposer replaced by a fixed hypothesis set (the existing concept classes and
    ``unknown'') and template-based missions, the loop resolves fewer questions about objects and failures outside
    the built-in models, and acquires new categories more slowly.
    \item A compositional generator trained only on the robot's own experience (format without culture) recovers
    part of the advantage in P6 within a familiar environment, and less of it the further a situation departs from
    what the robot has seen.
\end{enumerate}
A failure of P1 would indicate that the epistemic value of language-mediated questions does not exceed that of
the existing exploration; a failure of P4 would indicate that the loop does not transfer competence to the
automatic layer, which is its principal claim; and a failure of P6 would indicate that the language model adds
nothing a fixed hypothesis set does not, which would remove the case for including it. P7 separates the
contribution of language as a format for thought from that of the knowledge it carries.

\subsection{Limitations}
\label{sec:inner_speech:limitations}

The demand signal \eqref{eq:inner_demand} sums heterogeneous prediction errors whose precisions must themselves be
calibrated, and a poorly scaled term will either suppress the voice or keep it permanently engaged. A single
language model serves the person, the focused loop and the wandering loop; pre-emption keeps the person first but
bounds how much offline reasoning a busy robot can do. The tempered prior \eqref{eq:inner_prior} is a heuristic
combination rather than a derived posterior. The generator of P7 is substantially harder to build than the
ablation of P6, and may have to be approximated by a small model trained on the robot's logs and mission corpus,
whose pre-training would leak some culture back in. Finally, the analogy with human inner speech is functional, not
mechanistic: the proposal is that the same computational role --- serial, compositional self-questioning in the
service of epistemic value --- is useful for a robot, not that the robot's voice resembles the human one in any
further respect.
